\ficha{Flandreau e Flores (2012), The Peaceful Conspiracy}{flandreauflores2012io}

\begin{identificacao}
FLANDREAU, Marc; FLORES, Juan H. The Peaceful Conspiracy: Bond Markets and
International Relations During the Pax Britannica. \textit{International
Organization}, v. 66, n. 2, p. 211-241, 2012. DOI: 10.1017/S0020818312000070.

Artigo de periódico, 31 páginas, em inglês, lido integralmente na versão
publicada.
\end{identificacao}

\bloco{Problema e tese}

Polanyi atribuiu os cem anos de paz entre 1815 e 1914 à ação de um grupo de
banqueiros, a Haute Finance, mas postulou o mecanismo em vez de demonstrá-lo.
Os autores perguntam se o mercado de capitais de Londres discriminava a guerra
e por quê. O poder de monopólio existia, mas não era cartel: vinha do
prestígio como ativo proprietário de certificação. A condicionalidade pela paz
existia, mas não era bem público: as casas de topo barravam as guerras que
ameaçavam a própria marca e financiavam as que não ameaçavam.

\bloco{Argumento}

\begin{enumerate}
\item A certificação tende ao monopólio natural. Quem vende selo de qualidade
ganha ao mentir uma vez e perde receita futura ao ser desmascarado, e a
veracidade só é equilíbrio quando a participação de mercado já é grande o
bastante para que a perda supere o ganho (p. 215).

\item Os dados sustentam a premissa. Os Rothschild lideram as tabelas de
emissão de dívida soberana em Londres em todo o período, os três maiores
originadores sempre somam mais de 50\% do mercado e o índice de
Herfindahl-Hirschman vai de 1.196 a 2.432 conforme a fase (p. 221). Prestígio
é custo afundado e barreira à entrada, não acordo entre participantes, e por
isso escapa do ceticismo olsoniano quanto à sobrevivência de cartéis (p. 222).

\item O selo tem preço. Na entrada do Brasil em Londres, em 1825, a emissão
por uma casa comum fracassou e o anúncio de que os Rothschild tomariam o saldo
derrubou o spread brasileiro entre 50 e 100 pontos base, sem notícia nova
sobre o país (p. 224). Entre 1845 e 1877, novos devedores sem mácula pagavam
cerca de 300 pontos base a menos pelos Rothschild (p. 226), e trocas de
underwriter movem o spread em 127 a 150 pontos base para países limpos e cerca
de 400 no caso espanhol (p. 229).

\item O enforcement é delegado, e é isso que o torna crível. Colocada a
emissão, quem carrega a perda do calote é o investidor, não o banqueiro, cuja
preocupação é certificar no futuro. Punir o caloteiro passa a ser polir o
próprio escudo, o que resolve a inconsistência temporal que trava credores
dispersos com dinheiro preso no devedor (p. 222).

\item A condicionalidade aparece em contrato. O artigo 9 do primeiro
empréstimo internacional da Bélgica, de dezembro de 1831, dava aos
underwriters o direito de suspender os desembolsos caso uma potência
interviesse fora do mandato da Conferência de Londres (p. 230). O
parcelamento tornava a ameaça monitorável, e a recusa aberta de 1839 levou o
governo belga a assinar o tratado de paz (p. 233).

\item A paz, porém, não era bem público. Entre 1845 e 1913, 19,6\% das guerras
obtiveram financiamento em Londres (p. 218). Guerras entre potências
praticamente nunca, e as exceções das casas de topo são conflitos de desfecho
previsível e sem consequência geopolítica, entre eles os empréstimos dos
Rothschild ao Brasil em 1852 e 1865 (p. 235).
\end{enumerate}

\chave{Haute Finance was not the Red Cross}{p. 235}

\bloco{Conceitos e definições operacionais}

Prestígio é participação de mercado grande e persistente que encarece a
mentira, medida por tabelas de emissão, capital e desempenho das emissões (p.
222). Certificação é o serviço pelo qual o nome do intermediário substitui a
informação que o investidor não tem sobre o tipo do devedor, e seu valor
cresce com a assimetria (p. 226). Conditional lending é a subordinação
contratual do desembolso à política adotada pelo devedor (p. 230). Seasoning,
tomado de Tomz, distingue devedor novo de veterano e limpo de manchado (p.
226). O artigo não define padrão-ouro, conversibilidade nem lastro: seu objeto
é a estrutura do mercado de dívida, não o regime monetário.

\bloco{Evidência e método}

Base própria de emissões estrangeiras em Londres entre 1815 e 1913, cruzada
com a lista de guerras de Gleditsch, revisão do Correlates of War, e com
Reiter e Stam, organizada em quatro fases de auge e colapso do crédito externo
entre 1818 e 1913. São 380 empréstimos, dos quais 19 em tempo de guerra, 5\%
do total (p. 218). O aparato é descritivo: médias de prêmio de risco no
lançamento contra o consol britânico, tabulação por tipo de guerra e
comparação de pares antes e depois de troca de underwriter. Soma-se a
pesquisa nos contratos originais do arquivo Rothschild em Roubaix e a leitura
do \textit{Times} de Londres. O contrafactual é o cão que não latiu, a guerra
belga que não houve.

\bloco{Agentes e vertentes}

Polanyi é o alvo, e antes dele a tradição conspirativa de Hobson, Lysis,
Hilferding e Lenin, que via a mesma centralidade dos Rothschild com sinal
invertido. Do lado adversário estão Olson, Schultz e Weingast, céticos quanto
à coordenação de credores, Kirshner, para quem o financista é difuso e
apaziguador, e Tomz, para quem os credores são atomizados e a reputação do
país se acumula por bom comportamento. As casas: Rothschild e Baring no topo,
com o Baring como ameaça de entrada que disciplinava o Rothschild, e Morgan,
HSBC, Erlanger e Schroders como casas comuns.

\bloco{Críticas e limites}

Quinze a dezenove empréstimos de guerra, sem inferência estatística formal nem
controle multivariado, apenas médias de célula. A amostra é só Londres, com
Paris e Amsterdã excluídos por decisão explícita (p. 218), o que tende a
superestimar o embargo. A recusa de uma casa de topo é observacionalmente
equivalente a má qualidade do devedor, e a endogeneidade é tratada por leitura
de arquivo, não por desenho. O recorte é guerra, não política monetária: nada
no artigo diz que os mesmos bancos impusessem conversibilidade, paridade ou
disciplina fiscal. A cronologia também limita. O núcleo do enforcement
documentado é 1830-1870, e o período 1896-1913, o do Brasil da Caixa de
Conversão, é o de menor concentração da série, com índice de 1.196, o que
sugere poder de monopólio em queda às vésperas de 1906. O mecanismo de
condicionalidade financeira está desenvolvido no artigo dos mesmos autores na
\textit{European Review of Economic History}, fichado sob
\texttt{flandreauflores2012}, e não se repete aqui.

\begin{dialogo}
\item Procurar o argumento de que a Caixa de Conversão era exigência ou
expectativa dos credores externos, com o vocabulário ``crédito no exterior'',
``praça de Londres'', ``nossos banqueiros''. Se a defesa ortodoxa se apoia no
aval externo, o mecanismo em jogo é o enforcement por intermediário
certificador, e não o compromisso institucional interno.
\item Procurar menção nominal aos Rothschild fora do noticiário de cotação,
sobretudo em torno do funding loan de 1898 e da criação da Caixa. Testa se a
casa credora aparece como árbitro da política monetária ou apenas como agente
financeiro.
\item Procurar como a imprensa trata a agência de Londres e a custódia do ouro
em 1911. Defesa do depósito em Londres como sinal aos credores é certificação
por procuração; ataque como submissão é a leitura conspirativa de Hobson e
Lysis, contemporânea do debate.
\item Procurar o vocabulário de personificação do mercado, ``oligarquia
financeira'', ``trust do dinheiro'', ``sindicato de banqueiros'', em Alcindo
Guanabara, Barbosa Lima e no Correio da Manhã. Lysis publicou em 1908, dentro
da janela, e a hipótese é que a crítica expansionista importe essa gramática.
\end{dialogo}

\usonocapitulo{Parte II, como apoio teórico. Sustenta que o enforcement das
regras ortodoxas na periferia não decorria de um mercado difuso nem de
instituição doméstica, e sim de intermediários com poder de monopólio capazes
de condicionar crédito a política. Serve também para qualificar essa
afirmação, porque o mecanismo está documentado para a guerra, não para o
regime monetário, e para o meio do século XIX, não para 1906. O artigo não
trata do Brasil da Caixa de Conversão, e suas menções ao país são de 1825,
1852 e 1865. Para o conditional lending, remeter à ficha
\texttt{flandreauflores2012}.}
